\section{Descent to shorter intervals}\label{sec:descent}

We now combine the pair cancellation estimate with the seminorm furnished
by the tuple laws.  The quantity to be decreased is a diagonal tuple
integral of a Fourier lift.  Restriction to progressions compares this
quantity with the same quantity on shorter intervals; avoidance makes a
fixed proportion of the residue assignments cancel, apart from the small
exceptional measures.  This is the recurrence method of
\cite{IntersectivePower}, adapted here to preserve the prime argument
condition throughout the auxiliary family.

All structural parameters, including $M_0,d,\gamma$ and those in
\eqref{eq:parameters}, remain fixed.  At every integer scale $N\ge1$ we use
$\mathcal P_N,Q,H$ from \eqref{eq:kernel-cutoffs}.  For $l\ge1$ and
$A\subseteq[N]$, define
\begin{equation}\label{eq:descent-quantity}
 \begin{split}
 Y_l(N,A)&=Z_{l,\mathcal P_N}
       \bigl(\mathcal L_{[N],\mathcal P_N,Q}\mathbf1_A\bigr),\\
 \mathcal Y_l(N)&=\max_{\substack{A\subseteq[N]\\A\text{ avoiding at }l}}
                         Y_l(N,A).
 \end{split}
\end{equation}
Here $Y_l$ is unrelated to the prime-sum range $Y$ used in the preceding
section.  Lemma~\ref{lem:seminorm} and the constant Fourier coefficient of
the lift give the first bound in
\begin{equation}\label{eq:descent-trivial}
 \left(\frac{|A|}{N}\right)^d\le Y_l(N,A),
 \qquad 0\le\mathcal Y_l(N)\ll N^{o(1)}.
\end{equation}
For the upper bound, ordinary H\"older with the uniform single-slot
marginals gives $Z_{l,\mathcal P_N}(g)\le\|g\|_d^d$.  Apply
Proposition~\ref{prop:lift-moments}; its length hypothesis holds for large
$N$ since $6\beta<1$.  For bounded $N$, bounding each term of the finite
Fourier sum gives the same uniform conclusion.  In particular, the upper
bound in \eqref{eq:descent-trivial} holds for every $l\ge1$, with no
restriction of the form $l\le N^\zeta$.

For each $s\bmod M_0$ put
\[
 f_s(x)=M_0\mathbf1_A(x)\mathbf1_{\{x\equiv s\pmod{M_0}\}},
 \qquad G_s=\mathcal L_{[N],\mathcal P_N,Q}f_s.
\]
Averaging $f_s$ over uniform $s$ gives $\mathbf1_A$.  Thus, for independent
uniform residues $s_v\bmod M_0$ indexed by $v\in V$, multilinearity gives
\begin{equation}\label{eq:residue-average}
 Y_l(N,A)=\mathbb E_{(s_v)_{v\in V}}
          \mathcal Z_{l,\mathcal P_N}((G_{s_v})_{v\in V}).
\end{equation}
We shall bound these integrals in two ways.  The comparison with shorter
intervals holds for every $l$; cancellation will require the analytic
range $l\le N^\zeta$.

\subsection{Specialization and progression coordinates}

Fixing prime coordinates of a lift restricts its underlying set to a
residue class, provided the remaining Fourier denominator is sufficiently
small.  The following exact formula records the normalization.

\begin{lemma}[Residue specialization]\label{lem:residue-specialization}
Let $N\ge1$, $A\subseteq[N]$, $s\bmod M_0$, and let $I\subseteq[N]$ be a
nonempty consecutive interval.  Let $B\subseteq\mathcal P_N$, write
$q_B=\prod_{p\in B}p$, $D=M_0q_B$, and
$\mathcal S=\mathcal P_N\setminus B$.  Fix the coordinates indexed by $B$
of $\mathcal L_{I,\mathcal P_N,Q}f_s$ at $b=(b_p)_{p\in B}$, obtaining
$g_I:X_{\mathcal S}\to\R$.  Let $a'\bmod D$ be the unique residue
satisfying $a'\equiv s\pmod{M_0}$ and $a'\equiv b_p\pmod p$ for $p\in B$.
Then every frequency $\xi$ on $X_{\mathcal S}$ with $q(\xi)q_B\le Q$
satisfies
\begin{equation}\label{eq:residue-specialization}
 \widehat{g_I}(\xi)=\frac{D}{|I|}
       \sum_{\substack{x\in I\\x\equiv a'\pmod D}}
                       \mathbf1_A(x)e(-x\xi).
\end{equation}
\end{lemma}

\begin{proof}
The coefficient on the left is the sum of
$\widehat{f_s}_I(\xi+\omega)e(b\omega)$ over the frequencies $\omega$ on
$X_B$ retained by the lift.  The denominator condition retains every
$\omega$, because $q(\xi+\omega)=q(\xi)q(\omega)\le Q$.  Orthogonality of
these $q_B$ characters imposes $x\equiv b_p\pmod p$ for every $p\in B$
and contributes the factor $q_B$.  Together with the factor $M_0$ in
$f_s$, this gives \eqref{eq:residue-specialization}.
\end{proof}

\begin{proposition}[Comparison with shorter intervals]\label{prop:shorter-scale}
For all sufficiently large $N$, every $l\ge1$, and every set $A\subseteq[N]$
avoiding at $l$, the following holds.  Let $B\subseteq\mathcal P_N$ satisfy
$q_B\le N^\tau$, and put
\[
 D=M_0q_B,\qquad L_D=\lambda(D),\qquad
 n_B=\left\lceil\frac{N}{L_D}\right\rceil,\qquad
 \Lambda_B=\frac{L_Dn_B}{N},\qquad
 \mathcal S=\mathcal P_N\setminus B.
\]
For each $v\in V$, choose $s_v\bmod M_0$ and fix the $B$-coordinates of
$G_{s_v}$ arbitrarily, obtaining $g_v:X_{\mathcal S}\to\R$.  Then
\begin{equation}\label{eq:shorter-scale}
 \bigl|\mathcal Z_{l,\mathcal S}((g_v)_{v\in V})\bigr|
 \le\Lambda_B^d\mathcal Y_{lD}(n_B)+O(N^{-\sigma}),
\end{equation}
uniformly in all these choices.  The new scale and normalization satisfy
\begin{equation}\label{eq:shorter-scale-size}
 N^{1/2}\le M_0^{-k}N^{1-k\tau}\le n_B<N,
 \qquad 1\le\Lambda_B\le1+M_0^kN^{k\tau-1}.
\end{equation}
\end{proposition}

\begin{proof}
The scale bounds follow from $D\le L_D\le D^k$ and $M_0\ge2$; their first
inequality holds for sufficiently large $N$ because $k\tau<1/2$.
In particular, with $Q'=n_B^\beta$,
\begin{equation}\label{eq:shorter-cutoffs}
 H\le Q'\le Q.
\end{equation}
By Proposition~\ref{prop:signed-holder}, it is enough to prove
\eqref{eq:shorter-scale} for $Z_{l,\mathcal S}(g)$, where $g$ is any one of
the specialized inputs.

First keep only the frequencies of $g$ with denominator at most $H$,
and call the resulting real function $g_{\le H}$.  Telescope the diagonal
integral one slot at a time.  Lemma~\ref{lem:large-denominator}, followed
by Proposition~\ref{prop:lift-moments} and
Lemma~\ref{lem:lift-specialization}, bounds the change by
\begin{equation}\label{eq:specialization-tail}
 \bigl|Z_{l,\mathcal S}(g)-Z_{l,\mathcal S}(g_{\le H})\bigr|
 \ll H^{-\gamma}q_BN^{o(1)}=O(N^{-\sigma}).
\end{equation}
Indeed specialization costs $q_B^{1/2}$ in the $L^2$ norm of the removed
part, and $q_B^{1/(2(d-1))}$ in each of the other $d-1$ norms.  Their
product is $q_B$.  All frequency selections retain whole exact supports
before specialization, as required by the moment bound.
Here $N\ge10Q^6$ for large $N$, and the saving follows from
$\beta\gamma/2-\tau>\sigma$.

The auxiliary-family identity transports avoidance along progressions
of step $L_D=\lambda(D)$.  Since $D\mid L_D$, we split the prescribed
class modulo $D$ into its $L_D/D$ classes modulo $L_D$.
Let $a'\bmod D$ be the class determined by this specialization.  For
$j\in[L_D]$ with $j\equiv a'\pmod D$, set
\[
 c_j=j-L_D,\qquad
 A'_j=\{x'\in[n_B]:L_Dx'+c_j\in A\},\qquad
 \mathcal H_j=\mathcal L_{[n_B],\mathcal S,H}\mathbf1_{A'_j}.
\]
Every integer of $A$ in the class $a'\bmod D$ has a unique representation
$L_Dx'+c_j$ of this form: take $j$ to be its representative in $[L_D]$.
The definition by membership in $A$ excludes any surplus points beyond
$N$ in these progressions.

Each $A'_j$ avoids at $lD$.  More explicitly, a nonzero difference
$h_{lD}(m)$ with $m\ge1$ and $T_{lD}(m)$ prime would give the difference
$\lambda(D)h_{lD}(m)=h_l(Dm-i)$ in $A$, where $0\le i<D$ is supplied by
\eqref{eq:auxiliary-change}.  We have $Dm-i\ge1$ and
$T_l(Dm-i)=T_{lD}(m)$, so this would violate avoidance at $l$.
The nonzero condition is preserved because $\lambda(D)>0$.

Since $q_BH\le Q$, Lemma~\ref{lem:residue-specialization} now gives the
exact identity
\begin{equation}\label{eq:progression-fibers}
 g_{\le H}(z)=\frac{\Lambda_B}{L_D/D}
       \sum_{\substack{j\in[L_D]\\j\equiv a'\pmod D}}
       \mathcal H_j\bigl(L_D^{-1}(z-c_j)\bigr).
\end{equation}
To check it, every prime divisor of $L_D$ divides $D$, so multiplication
by $L_D$ is invertible on $X_{\mathcal S}$ and preserves Fourier
denominators.  For $q(\xi)\le H$, the coefficient of the $j$th summand at $\xi$ is
$\widehat{\mathcal H_j}(L_D\xi)e(-c_j\xi)$.  Multiplication by the prefactor
$\Lambda_B/(L_D/D)=Dn_B/N$ and summation over $j$ therefore gives
$(D/N)\sum_{x\in A,\ x\equiv a'\ (D)}e(-x\xi)$, which is exactly the
coefficient in \eqref{eq:residue-specialization} for $I=[N]$.
Both sides of \eqref{eq:progression-fibers} have zero coefficients at
denominators greater than $H$.

There are $L_D/D$ summands.  The triangle inequality from
Lemma~\ref{lem:seminorm} and its transport identity
\eqref{eq:lift-transport} imply
\begin{equation}\label{eq:fiber-seminorm}
 Z_{l,\mathcal S}(g_{\le H})
       \le\Lambda_B^d\max_j Z_{lD,\mathcal S}(\mathcal H_j).
\end{equation}
Increase the cutoff in $\mathcal H_j$ from $H$ to $Q'$.  The same
large-denominator argument, now without specialization, changes its
diagonal integral by $O(H^{-\gamma}N^{o(1)})=O(N^{-\sigma})$.
The moment hypothesis is valid at the new scale:
$n_B\ge10(Q')^6$ for large $N$, by $6\beta<1$ and
\eqref{eq:shorter-scale-size}.

The extended lift is obtained from
$\mathcal L_{[n_B],\mathcal P_{n_B},Q'}\mathbf1_{A'_j}$ by averaging out
its coordinates in $B\cap\mathcal P_{n_B}$ and adding unused coordinates
in $\mathcal S\setminus\mathcal P_{n_B}$.  Averaging contracts the
seminorm, and adding unused coordinates does not change it, by
Lemma~\ref{lem:seminorm}.  Consequently its diagonal integral is at most
$Y_{lD}(n_B,A'_j)\le\mathcal Y_{lD}(n_B)$.
Combining this observation with \eqref{eq:specialization-tail} and
\eqref{eq:fiber-seminorm} proves the diagonal bound, hence the stated
multilinear bound.  None of these steps restricts the size of $l$.
\end{proof}

\subsection{Cancellation for admissible residue assignments}

The preceding proposition controls every assignment in
\eqref{eq:residue-average}.  We now obtain a stronger bound for those
assignments whose cycle increments can occur as local polynomial values.
Call $(s_v)_{v\in V}$ \emph{good at $l$} if every
$s_{v_{j+1}}-s_{v_j}$, $0\le j<L$, is admissible modulo $M_0$ at $l$.
Let $\rho_l$ be the proportion of good assignments.  The zero-sum
admissible sequences of Lemma~\ref{lem:admissible-cycle} give
\begin{equation}\label{eq:good-proportion}
 \rho_l\ge\rho_*:=M_0^{1-L}>0.
\end{equation}
Indeed combine such a sequence at each prime power dividing $M_0$ by the
Chinese remainder theorem.  After prescribing these $L$ increments, the
initial cycle residue and every residue off the cycle are free.

\begin{proposition}[Cancellation on the nonexceptional laws]
\label{prop:descent-cancellation}
For all sufficiently large $N$, all $1\le l\le N^\zeta$, every
$A\subseteq[N]$ avoiding at $l$, and every good assignment $(s_v)_{v\in V}$,
\begin{equation}\label{eq:descent-cancellation}
 \left|\int\prod_{v\in V}G_{s_v}(z_v)\,
       d\bigotimes_{p\in\mathcal P_N}(\mu_{l,p}-\varepsilon_{l,p})\right|
       \ll N^{-\sigma}.
\end{equation}
The implied constant is uniform in $l$, $A$, and the assignment.
\end{proposition}

\begin{proof}
The aim is to leave just the two endpoint lifts on an ordered cycle
edge.  We first average over short intervals, then truncate and expand
the other inputs.  Each expansion term uses prime coordinates with
small product; conditioning on those coordinates leaves the pair to
which Proposition~\ref{prop:pair-cancellation} applies.

It suffices to use the probability laws $\mu_{l,p}^\circ$, since restoring
the omitted masses multiplies the integral by
$\prod_p(1-\eta_{l,p})\le1$.
Partition $[N]$ into $m=\lfloor N^\tau\rfloor$ consecutive intervals
of lengths as equal as possible.  Each $G_s$ is the weighted average of
$\mathcal L_{I,\mathcal P_N,Q}f_s$ over these intervals, with weights
$|I|/N$.  Thus the integral is an average over independently chosen
intervals for all slots.  Their lengths are comparable to $N^{1-\tau}$,
so the lift moment bound applies because $6\beta+\tau<1$.
Ordinary H\"older and uniform marginals bound every interval assignment
by $N^{o(1)}$ in absolute value.

The probability that all $L$ cycle slots use the same interval is
$\sum_I(|I|/N)^L=O(N^{-(L-1)\tau})$.  These assignments contribute
$O(N^{-\sigma})$.  For any remaining assignment, the cyclic sequence of
interval indices is nonconstant, and hence has a strict increase along
some directed cycle edge.  Choose such an edge, with first endpoint in
$I_1$ and second endpoint in the later interval $I_2$.

Keep the two endpoint lifts unchanged and truncate every other lift to
denominators at most $T=N^\tau$.  Telescoping with
Lemma~\ref{lem:large-denominator} and the lift moment bound costs
$O(T^{-\gamma}N^{o(1)})=O(N^{-\sigma})$.  Expand the truncated $d-2$
inputs into characters.  Their coefficients have absolute value at most
$M_0$, and each input has $O(T^2)$ frequencies, so the sum of absolute
coefficients in this product expansion is $O(T^{2(d-2)})$.

Fix one term.  Let $B$ be the union of the prime supports of its
characters and write
\[
 R=q_B\le T^{d-2},\qquad D=M_0R,\qquad
 \mathcal S=\mathcal P_N\setminus B.
\]
Condition on the entire tuple of labels at every prime in $B$, taking
any conditioning with positive probability under the product of the
$\mu_{l,p}^\circ$.  The expanded characters become constants of modulus
one.  The specialized endpoint lifts give functions $g_1,g_2$ on
$X_{\mathcal S}$, and the laws at primes in $\mathcal S$ remain unchanged.
Translation invariance shows that their pair integral is
\begin{equation}\label{eq:conditioned-pair}
 \sum_{\xi\text{ on }X_{\mathcal S}}
       \widehat g_1(-\xi)\widehat g_2(\xi)\Gamma(\xi),
\end{equation}
where $\Gamma$ is the multiplier for this directed cycle edge.  Indeed
only opposite characters survive common translation, and their product
is the character of the second label minus the first.
By \eqref{eq:edge-multiplier-bound}, Parseval, and the specialized moment
bound, the part with $q(\xi)>H$ is at most
\begin{equation}\label{eq:conditioned-pair-tail}
 H^{-\gamma}\|g_1\|_2\|g_2\|_2
       \ll H^{-\gamma}R N^{o(1)}.
\end{equation}

For the remaining frequencies, $HR\le Q$ and $D\le N^\beta$ for large
$N$, since $(d-2)\tau<\beta/2$.  Apply
Lemma~\ref{lem:residue-specialization} to the endpoint lifts, obtaining
classes $a'_1,a'_2\bmod D$.  Their directed difference is admissible at
$l$: goodness supplies admissibility modulo $M_0$, and the cycle
property of $\mu_{l,p}^\circ$ supplies it at the conditioned primes
$p\in B$.
Set
\[
 J_i=\mathbf1_A\mathbf1_{I_i}\mathbf1_{\{x\equiv a'_i\pmod D\}},
 \qquad i=1,2,
\]
as functions on $[N]$.  Every contributing pair $x,y$ has $x<y$ because
$I_1$ precedes $I_2$.  Avoidance therefore gives
$y-x\notin\mathcal D_l$.  All hypotheses of
Proposition~\ref{prop:pair-cancellation} hold.
Its frequency set $\Xi$ is exactly the remaining frequencies on
$X_{\mathcal S}$ with denominator at most $H$: coprimality with $D$
excludes the primes below $P$ and the primes in $B$, while every prime
in a denominator at most $H$ is at most $Q$.
The coefficient identity becomes
\[
 \widehat g_i(\xi)=\frac{DN}{|I_i|}\widehat{J_i}_{[N]}(\xi)
 \qquad(q(\xi)\le H).
\]
Consequently \eqref{eq:pair-cancellation} bounds the low-denominator
part of \eqref{eq:conditioned-pair} by
\begin{equation}\label{eq:conditioned-pair-main}
 O\left(H^{-\gamma/2}D\frac{N^2}{|I_1||I_2|}\right)
       =O(H^{-\gamma/2}DN^{2\tau}).
\end{equation}

After multiplication by the expansion bound $O(T^{2(d-2)})$, both
\eqref{eq:conditioned-pair-tail} and \eqref{eq:conditioned-pair-main}
are bounded by
\[
 N^{-\beta\gamma/4+(3d+2)\tau+o(1)}=O(N^{-\sigma}),
\]
since \eqref{eq:parameters} gives
$\beta\gamma/4>(3d+2)\tau+\sigma$.
These estimates are uniform in the conditioning, in the selected edge,
and in the interval assignment.  Averaging, and including the earlier
truncation and coincident-interval errors, proves the proposition.
\end{proof}

\subsection{The contracting recurrence}

We have obtained cancellation after removing the exceptional part of
every local law.  To return to the full laws in
\eqref{eq:residue-average}, we expand that removal.  The total mass of
the resulting exceptional terms is small enough to pay for their
possibly shorter progression scales.
For $B\subseteq\mathcal P_N$ write
\[
 \eta_{l,B}=\prod_{p\in B}\eta_{l,p},\qquad
 D_B=M_0q_B,\qquad
 n_B=\left\lceil\frac{N}{\lambda(D_B)}\right\rceil.
\]
Thus $D_\varnothing=M_0$.
The mass estimate in Proposition~\ref{prop:tuple-laws} and the fixed
choice of $P$ give, uniformly in $l$,
\begin{equation}\label{eq:exceptional-budget}
 \sum_{\varnothing\ne B\subseteq\mathcal P_N}q_B\eta_{l,B}
 \le\prod_{p\ge P}(1+p^{-3})-1\le\frac14.
\end{equation}

\begin{proposition}[Recurrence]\label{prop:descent-recurrence}
There is a constant $C_1$ such that for every sufficiently large $N$ and
$1\le l\le N^\zeta$,
\begin{equation}\label{eq:descent-recurrence}
 \begin{split}
 \mathcal Y_l(N)\le {}&C_1N^{-\sigma}+(1+u_N)\biggl[
 (1-\rho_l)\mathcal Y_{lM_0}(n_\varnothing)\\
 &\hspace{12mm}+\rho_l
 \sum_{\substack{\varnothing\ne B\subseteq\mathcal P_N\\q_B\le N^\tau}}
       \eta_{l,B}\mathcal Y_{lD_B}(n_B)\biggr],\\
 u_N={}&(1+M_0^kN^{k\tau-1})^d-1=o(1).
 \end{split}
\end{equation}
The constant and the sufficiently-large threshold are independent of $l$.
\end{proposition}

\begin{proof}
Fix an avoiding set $A$.  For assignments in \eqref{eq:residue-average}
that are not good, apply Proposition~\ref{prop:shorter-scale} with
$B=\varnothing$.  Their proportion is $1-\rho_l$.
For a good assignment, Proposition~\ref{prop:descent-cancellation}
handles the product of the measures $\mu_{l,p}-\varepsilon_{l,p}$.
The difference from the full product is
\[
 \bigotimes_{p\in\mathcal P_N}\mu_{l,p}
 -\bigotimes_{p\in\mathcal P_N}(\mu_{l,p}-\varepsilon_{l,p})
 =\sum_{\varnothing\ne B\subseteq\mathcal P_N}(-1)^{|B|+1}
   \left(\bigotimes_{p\in B}\varepsilon_{l,p}\right)
   \otimes\left(\bigotimes_{p\notin B}\mu_{l,p}\right).
\]
The positive product measure in the $B$th summand has mass
$\eta_{l,B}$.  If this mass is zero the term vanishes.  Otherwise,
after normalization, every single-slot marginal is uniform by
simultaneous translation invariance.  H\"older and
Proposition~\ref{prop:lift-moments} thus bound its integral by
$N^{o(1)}$ in absolute value.  The terms with $q_B>N^\tau$ have total
mass at most $N^{-\tau}/4$ by \eqref{eq:exceptional-budget}, and hence
contribute $O(N^{-\sigma})$.

For $q_B\le N^\tau$, condition on all labels at the primes in $B$.
The remaining law is the product of the full measures $\mu_{l,p}$, so
Proposition~\ref{prop:shorter-scale} applies to the specialized inputs.
Its factor $\Lambda_B^d$ is at most $1+u_N$ by
\eqref{eq:shorter-scale-size}, and its errors sum to $O(N^{-\sigma})$
by \eqref{eq:exceptional-budget}.  Taking absolute upper bounds in the
signed expansion and averaging over the good assignments gives their
contribution in \eqref{eq:descent-recurrence}.  Finally maximize over
$A$ in \eqref{eq:residue-average}.
\end{proof}

\subsection{Uniform induction and the density bound}

The recurrence uses prime estimates only when $l\le N^\zeta$, but its
successors $lD_B$ need not lie in the corresponding analytic range at
$n_B$.  We therefore prove a bound simultaneously for every auxiliary
parameter $l$.  The factor $l^{b_0}$ below covers the complementary
range by the elementary bound \eqref{eq:descent-trivial}.

\begin{proposition}[Uniform power decay]\label{prop:uniform-descent}
There are constants $a_0,b_0>0$ and $C_0\ge1$, depending only on $h$, such
that for all integers $l,N\ge1$,
\begin{equation}\label{eq:uniform-descent}
 \mathcal Y_l(N)\le C_0l^{b_0}N^{-a_0}.
\end{equation}
\end{proposition}

\begin{proof}
After all parameters in \eqref{eq:parameters} have been fixed, choose
$a_0>0$ sufficiently small and put $b_0=2a_0/\zeta$, so that
\begin{equation}\label{eq:induction-parameters}
 a_0<\sigma,\qquad s_0:=ka_0+b_0\le1,\qquad
 \theta_0:=M_0^{s_0}(1-3\rho_*/4)<1.
\end{equation}
This is possible because $s_0\to0$ with $a_0$ while $\rho_*>0$ is fixed.
If $l>N^\zeta$, then
\[
 l^{b_0}N^{-a_0}\ge N^{a_0}.
\]
The estimate \eqref{eq:descent-trivial}, used with exponent $a_0$, proves
\eqref{eq:uniform-descent} in this range after a fixed enlargement of
$C_0$.  This treats all such $l,N$ before the induction.

For the remaining pairs, induct on the integer $N$, simultaneously for
all $l\ge1$.  At a sufficiently large scale with $l\le N^\zeta$, every
successor in \eqref{eq:descent-recurrence} satisfies $n_B<N$ by
\eqref{eq:shorter-scale-size}.  The induction hypothesis is available at
$lD_B$ whether or not that parameter satisfies $lD_B\le n_B^\zeta$.
Since $n_B\ge N/D_B^k$, it gives
\[
 \mathcal Y_{lD_B}(n_B)
 \le C_0(lD_B)^{b_0}n_B^{-a_0}
 \le C_0l^{b_0}N^{-a_0}D_B^{s_0}.
\]
Use $q_B^{s_0}\le q_B$ and \eqref{eq:exceptional-budget} to bound the
bracket in \eqref{eq:descent-recurrence} by
\[
 \begin{split}
 C_0l^{b_0}N^{-a_0}M_0^{s_0}
   \bigl(1-\rho_l+\rho_l/4\bigr)
 &\le C_0l^{b_0}N^{-a_0}\theta_0,
 \end{split}
\]
where \eqref{eq:good-proportion} was used in the last step.
Choose a fixed threshold $N_*$, large enough for the recurrence and for
\[
 C_1N^{a_0-\sigma}+(1+u_N)\theta_0\le1
 \qquad(N\ge N_*).
\]
Such a threshold exists by \eqref{eq:induction-parameters} and $u_N=o(1)$;
it is independent of $C_0\ge1$.  Dividing the recurrence by
$C_0l^{b_0}N^{-a_0}$ now proves the induction step, since
$C_0l^{b_0}\ge1$.
Finally choose $C_0$ large enough to cover both the previously treated
range $l>N^\zeta$ and every $N<N_*$, uniformly in $l$, by
\eqref{eq:descent-trivial}.  This completes the induction.
\end{proof}

\begin{proof}[Proof of Theorem~\ref{thm:main}]
At $l=1$ we have $T_1(m)=m$ and $h_1=h$, so the hypothesis of the theorem
is exactly avoidance at $l=1$.  Combining \eqref{eq:descent-trivial}
with \eqref{eq:uniform-descent} gives
\[
 \left(\frac{|A|}{N}\right)^d\le C_0N^{-a_0},
 \qquad |A|\le C_0^{1/d}N^{1-a_0/d}.
\]
Thus one may take $c_h=a_0/d$ and $C_h=C_0^{1/d}$, for every $N\ge1$.
Only strictly ordered pairs, and hence positive differences, were used
in cancellation.  Restriction to progressions preserved nonzero
forbidden values and their prime arguments.  A value $h(p)=0$ therefore
imposes no additional condition at any stage.
\end{proof}
